\documentclass[11pt,a4paper]{article}

\usepackage[utf8]{inputenc}
\usepackage[T1]{fontenc}
\usepackage[spanish]{babel}
\usepackage{lmodern}
\usepackage{amsmath,amssymb}
\usepackage[a4paper,margin=1.5cm]{geometry}
\usepackage{array,booktabs,longtable}
\usepackage{hyperref}
\usepackage{xurl}
\usepackage{microtype}

\hypersetup{
	colorlinks=true,
	linkcolor=black,
	urlcolor=black
}

\setcounter{secnumdepth}{0}
\setlength{\parindent}{0pt}
\setlength{\parskip}{6pt}
\setlength{\tabcolsep}{3pt}
\renewcommand{\arraystretch}{1.15}
\newcolumntype{L}[1]{>{\raggedright\arraybackslash}p{#1}}
\setlength{\LTleft}{0pt}
\setlength{\LTright}{0pt}
\sloppy

\title{\textbf{Tunjo Fest Soundpainting}\\
	\large Propuesta para la Casa de la Cultura de Facatativá}
\author{Soundpainting Colombia}
\date{}

\begin{document}
	
	\maketitle
	
	\section{1 Presentación general de la propuesta}
	
	\subsection{1.1 Nombre, descripción y destinatarios}
	
	Tunjo Fest Soundpainting es una propuesta de formación, ensamble y presentación escénica basada en Soundpainting, un lenguaje de señas para la composición en tiempo real.
	
	La práctica permite integrar música, danza, teatro, circo, literatura, artes plásticas, medios audiovisuales y otras disciplinas mediante improvisación dirigida. Un director o SoundPainter utiliza señas para orientar entradas, silencios, repeticiones, variaciones, desplazamientos, acciones escénicas y relaciones entre intérpretes.
	
	La propuesta está diseñada para articular a artistas invitados con participantes de las escuelas de formación de la Casa de la Cultura de Facatativá. El proceso culmina con una función pública en el marco de Tunjo Fest, con una duración estimada entre 60 y 90 minutos.
	
	Está dirigida a estudiantes de las escuelas de formación de la Casa de la Cultura, docentes y directores de procesos artísticos, artistas locales, jóvenes, familias, comunidad general de Facatativá y público asistente a Tunjo Fest.
	
	\subsection{1.2 Colectivo proponente y ensamble invitado}
	
	Colectivo proponente: Soundpainting Colombia.
	
	Ensamble invitado: Ensamble Liminal.
	
	Artistas base del Ensamble Liminal: 11.
	
	Artistas invitados de circo de Facatativá: 3.
	
	Total de intérpretes invitados para la función: 14.
	
	Formato: formación, ensamble interdisciplinar y función pública.
	
	El Ensamble Liminal es el grupo artístico invitado para la realización de la función principal y el acompañamiento del proceso de ensamble. Su conformación integra músicos, actores, bailarines, performers, SoundPainters y artistas invitados de circo.
	
	El ensamble trabaja a partir de la improvisación escénica, la composición en tiempo real, la escucha colectiva y la dirección mediante señas. Su participación permite articular una experiencia interdisciplinar con las escuelas de formación de la Casa de la Cultura.
	
	\subsection{1.3 Contexto de realización y referencias}
	
	La función principal se proyecta para el marco de Tunjo Fest. El lugar de presentación será la tarima principal de Tunjo Fest, en octubre, según disponibilidad y condiciones técnicas.
	
	Video de referencia:
	
	\url{https://www.youtube.com/watch?v=CsyuLBz3cNg}
	
	Referencia visual de la propuesta:
	
	\url{https://canva.link/3oth4ravfmedgae}
	
	\url{https://www.canva.com/design/DAG4suoYzf4/IvxMuhm1ZMFyDcmN4m4vkw/edit}
	
	\subsection{1.4 Objetivos}
	
	\textbf{Objetivo general}
	
	Desarrollar una experiencia de Soundpainting en Facatativá mediante formación, ensamble interdisciplinar y función pública, integrando al Ensamble Liminal, artistas invitados de circo de Facatativá y participantes de las escuelas de formación de la Casa de la Cultura.
	
	\textbf{Objetivos específicos}
	
	\begin{itemize}
		\item
		Introducir a las escuelas de formación de la Casa de la Cultura en los principios básicos del Soundpainting.
		\item
		Crear espacios de trabajo entre música, danza, teatro, circo, literatura, artes visuales y medios audiovisuales.
		\item
		Realizar sesiones de ensamble con todas las disciplinas, articulando artistas invitados y artistas locales.
		\item
		Identificar Topics, Paletas y Texturas que puedan alimentar la composición escénica final.
		\item
		Presentar una función pública de composición escénica en tiempo real.
		\item
		Visibilizar los procesos artísticos de la Casa de la Cultura.
		\item
		Fortalecer la programación escénica y cultural en el marco de Tunjo Fest.
	\end{itemize}
	
	\section{2 Justificación cultural, social y territorial}
	
	\subsection{2.1 Pertinencia del Soundpainting}
	
	El Soundpainting es un lenguaje de señas creado por Walter Thompson para la composición en tiempo real. Funciona como herramienta de mediación artística, creación espontánea, juego escénico e intercambio de saberes.
	
	Su metodología permite que artistas de distintas disciplinas trabajen bajo un mismo sistema de comunicación escénica mediante acciones dirigidas, improvisación y composición colectiva.
	
	El Soundpainting permite la integración colectiva de distintos actores culturales, favorece la accesibilidad y democratización del arte, fortalece la seguridad expresiva de los participantes y promueve una escucha abierta que permite el diálogo creativo entre la diferencia. Esta orientación dialoga con experiencias como Casa Teatro Aguilar, gestionada por la Fundación Cultural U.I.I.L.H., con quienes se diseña esta propuesta para seguir abriendo espacios para gestores culturales, sociales y comunitarios que trabajan y transforman a la comunidad a través de las artes.
	
	\subsection{2.2 Pertinencia para Facatativá y la Casa de la Cultura}
	
	La Casa de la Cultura de Facatativá cuenta con procesos de formación en música, danza, teatro, literatura, artes visuales, medios audiovisuales y otras prácticas artísticas. Esta diversidad permite desarrollar una propuesta interdisciplinar que conecte a estudiantes, docentes, artistas invitados y comunidad en general.
	
	La propuesta permite que las escuelas de formación participen en una experiencia común sin abandonar su lenguaje propio. Cada área aporta sus herramientas, técnicas y materiales expresivos dentro de una estructura de creación compartida.
	
	\subsection{2.3 Aporte cultural, comunitario y relación con la agenda municipal}
	
	El proyecto fortalece el tejido cultural local al conectar artistas invitados, docentes, estudiantes, artistas locales y comunidad en una experiencia común de creación. Su pertinencia está en que articula procesos ya existentes de la Casa de la Cultura mediante una metodología compartida de composición escénica, sin necesidad de crear una escuela nueva.
	
	En el marco de Tunjo Fest, la propuesta funciona como una actividad artística, formativa y comunitaria que permite presentar un resultado público con participación local, acompañamiento de artistas profesionales e integración de lenguajes diversos.
	
	A través de la improvisación, la composición en vivo y la escucha colectiva, el proyecto favorece el diálogo entre escuelas artísticas, amplía las posibilidades escénicas del evento mediante la inclusión del circo y la integración potencial de medios audiovisuales, artes visuales, pintura y dibujo, y contribuye a visibilizar los procesos formativos de la Casa de la Cultura dentro de la agenda cultural municipal.
	
	\section{3 Diseño artístico, pedagógico y metodológico}
	
	\subsection{3.1 Descripción de la presentación}
	
	La función principal será una presentación escénica de Soundpainting construida en tiempo real. Los SoundPainters dirigirán a los intérpretes mediante señas, generando composiciones instantáneas a partir del sonido, el cuerpo, la voz, el movimiento, el texto, la imagen y la acción física.
	
	La muestra final se proyecta como una composición organizada en cuatro momentos, cada uno dirigido por dos SoundPainters. Esta estructura permitirá alternar perspectivas de dirección, ampliar las posibilidades de diálogo escénico y generar distintas relaciones entre los intérpretes.
	
	Cada momento podrá integrar materiales descubiertos durante el proceso de ensamble, especialmente Topics, Paletas y Texturas. Estos materiales no funcionarán como una partitura cerrada, sino como recursos disponibles para la composición en tiempo real.
	
	La presentación podrá articular cruces entre música, teatro, danza, circo, literatura, voz, artes plásticas, imagen en vivo, medios audiovisuales, registro y proyección. Estos lenguajes podrán combinarse tanto en composiciones específicas como en el ensamble general dirigido por los SoundPainters.
	
	\begin{itemize}
		\item
		Primer momento: apertura del Ensamble Liminal y presentación del lenguaje escénico.
		\item
		Segundo momento: composiciones por grupos pequeños o cruces disciplinares.
		\item
		Tercer momento: integración de artistas locales seleccionados y materiales desarrollados durante el proceso formativo.
		\item
		Cuarto momento: ensamble final interdisciplinario con participación ampliada y cierre colectivo.
	\end{itemize}
	
	\subsection{3.2 Áreas y escuelas participantes}
	
	La propuesta puede trabajar con hasta 7 áreas o grupos de formación, según disponibilidad de la Casa de la Cultura, número de estudiantes, horarios, espacios, instrumentos y condiciones técnicas.
	
	{\small
		\begin{longtable}{@{}L{0.27\textwidth}L{0.30\textwidth}L{0.37\textwidth}@{}}
			\toprule
			\textbf{Área o grupo}
			& \textbf{Participación estimada}
			& \textbf{Observaciones}
			\\
			\midrule
			\endhead
			\bottomrule
			\endlastfoot
			Música 1 & Parte de hasta 100 participantes entre los tres grupos de música & Grupo por definir con la Casa de la Cultura. Su participación depende de instalaciones, instrumentos y logística. \\
			Música 2 & Parte de hasta 100 participantes entre los tres grupos de música & Puede articularse con agrupaciones como sinfónica, fiestera, cuerdas, coral, estudiantina, vallenato o Big Band. \\
			Música 3 & Parte de hasta 100 participantes entre los tres grupos de música & La selección final dependerá de la disponibilidad de profesores, estudiantes, instrumentos y espacios. \\
			Danza 1 & Parte de 30 participantes aproximados entre los dos grupos de danza & Grupo por definir con la Casa de la Cultura. \\
			Danza 2 & Parte de 30 participantes aproximados entre los dos grupos de danza & Puede incluir urbano, ritmos latinos, folclor, ballet o contemporáneo. \\
			Teatro y literatura & Aproximadamente 15 participantes & Trabajo con voz, texto, acción, presencia escénica, escritura, lectura e improvisación. \\
			Medios audiovisuales, pintura y dibujo & Por definir & Puede apoyar registro, proyección, imagen en vivo, composición gráfica, dibujo, pintura o recursos visuales. Requiere definir condiciones de proyección y materiales. \\
		\end{longtable}
	}
	
	La participación en la función final no requiere la presencia total de todos los estudiantes. Se seleccionarán representantes por área, de acuerdo con la modalidad escogida, el espacio disponible y las condiciones técnicas.
	
	\subsection{3.3 Modalidades de implementación}
	
	La propuesta contempla tres modalidades. Cada una ajusta el alcance formativo, el número de sesiones, el tiempo de ensamble, la participación local y el presupuesto.
	
	{\footnotesize
		\begin{longtable}{@{}L{0.18\textwidth}L{0.18\textwidth}L{0.18\textwidth}L{0.18\textwidth}L{0.18\textwidth}@{}}
			\toprule
			\textbf{Modalidad}
			& \textbf{Escala}
			& \textbf{Formación}
			& \textbf{Ensamble}
			& \textbf{Función}
			\\
			\midrule
			\endhead
			\bottomrule
			\endlastfoot
			Descubriendo el lenguaje & Pequeña & 1 sesión de 2 horas por área & 2 horas & 60 a 90 minutos \\
			Explorando nuevas posibilidades & Mediana & 2 sesiones de 2 horas por área & 2 horas & 60 a 90 minutos \\
			Creación y composición intuitiva & Grande & 4 sesiones de 2 horas por área & 2 sesiones de ensamble general, equivalentes a 4 horas en total & 60 a 90 minutos \\
		\end{longtable}
	}
	
	\textbf{Descubriendo el lenguaje}
	
	Modalidad recomendada para una programación corta o una primera aproximación al Soundpainting. Incluye una sesión de formación de dos horas por área, dos horas de ensamble, participación puntual de artistas locales y función pública con 14 intérpretes invitados.
	
	\textbf{Explorando nuevas posibilidades}
	
	Modalidad recomendada para una agenda cultural con componente pedagógico y participación local media. Incluye dos sesiones de formación de dos horas por área, dos horas de ensamble, participación de artistas locales seleccionados y función pública con 14 intérpretes invitados.
	
	\textbf{Creación y composición intuitiva}
	
	Modalidad recomendada para Tunjo Fest o para una programación municipal de mayor alcance. Incluye cuatro sesiones de formación de dos horas por área o disciplina, dos sesiones de ensamble general equivalentes a cuatro horas en total, participación de escuelas de formación de la Casa de la Cultura, integración con el Ensamble Liminal y artistas invitados, selección de participantes locales para la función y función pública con 14 intérpretes invitados.
	
	\subsection{3.4 Metodología de formación, ensamble y función pública}
	
	La propuesta se organiza en tres momentos: formación, ensamble y función pública.
	
	\textbf{Formación}
	
	Las sesiones de formación introducen los principios básicos del Soundpainting y su aplicación a diferentes lenguajes artísticos.
	
	Se trabajan elementos como atención, escucha, respuesta, repetición, variación, entrada y salida, pausa, gesto, movimiento, sonido, texto, imagen y acción física.
	
	Cada área trabaja desde su propio lenguaje. La danza aborda movimiento, presencia y composición espacial. La música trabaja sonido, ritmo, textura y escucha. Teatro y literatura trabajan voz, palabra, cuerpo, acción, lectura, escritura, narración e improvisación. Circo trabaja precisión, acción física y composición visual. Medios audiovisuales, pintura y dibujo pueden integrarse mediante imagen, proyección, registro, intervención gráfica o composición visual.
	
	\textbf{Ensamble}
	
	El ensamble reúne participantes seleccionados de diferentes áreas. Su propósito es construir una práctica común entre artistas invitados y artistas locales.
	
	Durante el ensamble se practica el lenguaje en un contexto multidisciplinar, estableciendo relaciones entre disciplinas, entradas, salidas, transiciones y momentos de participación con contenidos previamente diseñados dentro de la función pública. De esta manera, se construye una narrativa colectiva que responde tanto a la expresión propia como a los sentidos individuales de los intérpretes.
	
	Además de practicar el lenguaje, el ensamble funciona como un laboratorio de descubrimiento, selección y organización de materiales escénicos para la presentación final. En este proceso se identifican Topics, Paletas y Texturas.
	
	Los Topics son temáticas, detonantes o ejes de improvisación que orientan una acción escénica, sonora, corporal, textual o visual. Permiten activar una situación, una imagen, una pregunta, una atmósfera o una relación entre disciplinas. Pueden surgir de ejercicios de creación, conversaciones con los participantes, imágenes del territorio, dinámicas de grupo o necesidades expresivas de la obra.
	
	Las Paletas son pequeñas piezas previamente preparadas. Pueden incluir secuencias sonoras, frases de movimiento, acciones teatrales, fragmentos textuales, imágenes, gestos, atmósferas o estructuras breves. Estas piezas pueden ser musicales, corporales, teatrales, textuales, visuales o interdisciplinares, y luego pueden ser retomadas, transformadas, reorganizadas o intervenidas por los SoundPainters durante la función. No se conciben como fragmentos cerrados, sino como materiales flexibles para la composición en tiempo real.
	
	Las Texturas son momentos de improvisación encontrados durante el ensamble que dejan una sensación específica. Pueden estar asociadas a una cualidad sonora, corporal, espacial, visual, rítmica, emocional o narrativa, y tener una cualidad de tensión, calma, acumulación, suspensión, juego, contraste, densidad, fragilidad, celebración, caos, silencio o escucha. Estas texturas se reconocen, se depuran y se conservan como posibles materiales escénicos para la muestra final.
	
	El proceso de ensamble permite que la función no sea únicamente una improvisación abierta, sino una composición en tiempo real alimentada por materiales descubiertos en la práctica. Así, la obra conserva la espontaneidad del Soundpainting, pero integra una memoria colectiva construida durante el proceso formativo.
	
	\textbf{Función pública}
	
	La función pública presenta el resultado del proceso ante la comunidad. Se realiza con el Ensamble Liminal, los artistas invitados de circo y participantes seleccionados de las escuelas de formación.
	
	Cada momento de la muestra final contará con dos SoundPainters, quienes dirigirán la composición en tiempo real y activarán los materiales construidos durante el proceso de formación y ensamble.
	
	La estructura final se definirá de acuerdo con la modalidad seleccionada, el número de participantes, el espacio disponible y las condiciones técnicas.
	
	\subsection{3.5 Metodología para la construcción de horarios y la organización de espacios}
	
	Antes del inicio de las clases o talleres, se realizará una conversación de coordinación con la Casa de la Cultura, las escuelas de formación y sus profesores, con el fin de organizar las condiciones generales del proceso.
	
	En este espacio se definirán las áreas participantes, el número aproximado de estudiantes por área, la disponibilidad de salones y espacios de ensayo, los horarios posibles y las necesidades técnicas de cada disciplina.
	
	También se revisarán las condiciones de participación en la función final y los requerimientos específicos de proyección, sonido, instrumentos o materiales visuales. La programación deberá ajustarse a las dinámicas reales de la Casa de la Cultura, los docentes, los estudiantes y los espacios disponibles. Por esta razón, los horarios definitivos se construirán de común acuerdo, según la modalidad seleccionada y la capacidad operativa de la institución.
	
	\section{4 Producción, espacios y requerimientos técnicos}
	
	\subsection{4.1 Espacios de formación, ensayo y función}
	
	Los espacios de formación y ensayo se definirán en coordinación con la Casa de la Cultura, de acuerdo con la disponibilidad, las condiciones técnicas y las necesidades de cada disciplina.
	
	El proceso podrá desarrollarse en la Casa de la Cultura, el Auditorio Teatro Municipal, el Teatro Municipal, salones amplios para danza, teatro y circo, espacios con condiciones adecuadas de sonido para música, lugares con conexión eléctrica para audiovisuales o proyección, y espacios apropiados para dibujo, pintura u otras artes visuales, si estas áreas se integran al proyecto.
	
	La función principal se proyecta para la tarima principal de Tunjo Fest, en octubre, según disponibilidad y condiciones técnicas. Para su realización se requiere un escenario con sonido, iluminación, condiciones adecuadas de circulación, piso seguro y visibilidad suficiente para que las señas del SoundPainter puedan ser vistas claramente por los intérpretes.
	
	\subsection{4.2 Requerimientos básicos del Ensamble Liminal}
	
	{\small
		\begin{longtable}{@{}L{0.32\textwidth}L{0.62\textwidth}@{}}
			\toprule
			\textbf{Componente}
			& \textbf{Requerimiento}
			\\
			\midrule
			\endhead
			\bottomrule
			\endlastfoot
			Teclado & Salida estéreo L/R por TRS o caja directa. \\
			Guitarra & Amplificador de guitarra, micrófono dinámico con soporte y/o caja directa. \\
			Voz / looper & Micrófono dinámico con soporte y salida L/R por TRS para conexión de looper. \\
			Voz con efectos & Dos micrófonos dinámicos, uno con efectos y otro sin efectos. \\
			Grupo musical & Retorno o monitoreo. \\
			Actores & Micrófonos inalámbricos de diadema, si el espacio lo requiere. \\
			Bailarines y circo & Piso seguro, espacio despejado y condiciones de iluminación. \\
			SoundPainters & Iluminación frontal suficiente para que las señas sean visibles. \\
		\end{longtable}
	}
	
	El rider técnico definitivo se entregará o ajustará después de conocer el espacio de función y la modalidad seleccionada.
	
	\section{5 Presupuesto flexible y condiciones de implementación}
	
	\subsection{5.1 Criterios generales del presupuesto}
	
	El presupuesto de Tunjo Fest Soundpainting se plantea de manera flexible, con base en un rango de 4 a 7 áreas participantes y de acuerdo con la modalidad seleccionada, la disponibilidad de escuelas, horarios, espacios, instrumentos y condiciones técnicas.
	
	Para facilitar su lectura institucional, se presentan primero los valores consolidados por modalidad y, posteriormente, un anexo presupuestal con los costos unitarios, fórmulas de cálculo y desgloses específicos.
	
	El cálculo considera el número de áreas participantes, las sesiones de formación por área, las horas y sesiones de ensamble, los intérpretes en función, los viáticos y los posibles rubros adicionales, como cobertura audiovisual, artistas invitados especializados o requerimientos técnicos complementarios.
	
	Los valores consolidados incluyen honorarios y viáticos de formación, honorarios y viáticos de ensamble, honorarios de intérpretes en función, la participación de 11 artistas del Ensamble Liminal y 3 artistas invitados de circo de Facatativá.
	
	Salvo acuerdo específico, estos valores no incluyen sonido, iluminación, tarima, planta eléctrica, transporte intermunicipal, hospedaje, alimentación general, registro o edición audiovisual, diseño gráfico, seguridad, primeros auxilios, permisos, alquiler de espacios, imprevistos, equipos de proyección, pantallas para trabajo audiovisual o visual, ni otros requerimientos técnicos adicionales.
	
	La cobertura audiovisual, los artistas invitados especializados y los requerimientos técnicos complementarios podrán sumarse al presupuesto base si son aprobados por la Casa de la Cultura, la Alcaldía, la organización de Tunjo Fest o la entidad responsable de la contratación. Estos componentes podrán ser asumidos directamente por dichas entidades o incluirse como rubros adicionales, según la modalidad seleccionada y los acuerdos establecidos.
	
	La implementación de la propuesta queda sujeta a la disponibilidad de espacios, instrumentos, profesores, escuelas participantes, condiciones técnicas, agenda institucional y acuerdos previos con la Casa de la Cultura.
	
	\emph{Todas las cifras monetarias de las tablas presupuestales están expresadas en pesos colombianos COP.}
	
	\subsection{5.2 Presupuesto consolidado por modalidad}
	
	Los siguientes valores corresponden al presupuesto base estimado para cada modalidad, calculado según el número de áreas participantes.
	
	\textbf{Modalidad 1: Descubriendo el lenguaje}
	
	Escala: pequeña.
	
	Formación: 1 sesión de 2 horas por área.
	
	Ensamble: 2 horas en 1 sesión de ensamble.
	
	Función: 14 intérpretes.
	
	Duración de la función: 60 a 90 minutos.
	
	{\footnotesize
		\begin{longtable}{@{}L{0.18\textwidth}L{0.18\textwidth}L{0.18\textwidth}L{0.18\textwidth}L{0.18\textwidth}@{}}
			\toprule
			\textbf{Número de áreas}
			& \textbf{Formación con viáticos}
			& \textbf{Ensamble con viáticos}
			& \textbf{Función}
			& \textbf{Total consolidado}
			\\
			\midrule
			\endhead
			\bottomrule
			\endlastfoot
			4 áreas & \$1.360.000 COP & \$2.540.000 COP & \$8.400.000 COP & \$12.300.000 COP \\
			5 áreas & \$1.700.000 COP & \$2.540.000 COP & \$8.400.000 COP & \$12.640.000 COP \\
			6 áreas & \$2.040.000 COP & \$2.540.000 COP & \$8.400.000 COP & \$12.980.000 COP \\
			7 áreas & \$2.380.000 COP & \$2.540.000 COP & \$8.400.000 COP & \$13.320.000 COP \\
		\end{longtable}
	}
	
	\textbf{Modalidad 2: Explorando nuevas posibilidades}
	
	Escala: mediana.
	
	Formación: 2 sesiones de 2 horas por área.
	
	Ensamble: 2 horas en 1 sesión de ensamble.
	
	Función: 14 intérpretes.
	
	Duración de la función: 60 a 90 minutos.
	
	{\footnotesize
		\begin{longtable}{@{}L{0.18\textwidth}L{0.18\textwidth}L{0.18\textwidth}L{0.18\textwidth}L{0.18\textwidth}@{}}
			\toprule
			\textbf{Número de áreas}
			& \textbf{Formación con viáticos}
			& \textbf{Ensamble con viáticos}
			& \textbf{Función}
			& \textbf{Total consolidado}
			\\
			\midrule
			\endhead
			\bottomrule
			\endlastfoot
			4 áreas & \$2.720.000 COP & \$2.540.000 COP & \$8.400.000 COP & \$13.660.000 COP \\
			5 áreas & \$3.400.000 COP & \$2.540.000 COP & \$8.400.000 COP & \$14.340.000 COP \\
			6 áreas & \$4.080.000 COP & \$2.540.000 COP & \$8.400.000 COP & \$15.020.000 COP \\
			7 áreas & \$4.760.000 COP & \$2.540.000 COP & \$8.400.000 COP & \$15.700.000 COP \\
		\end{longtable}
	}
	
	\textbf{Modalidad 3: Creación y composición intuitiva}
	
	Escala: grande.
	
	Formación: 4 sesiones de 2 horas por área.
	
	Ensamble: 4 horas en 2 sesiones de ensamble general.
	
	Función: 14 intérpretes.
	
	Duración de la función: 60 a 90 minutos.
	
	{\footnotesize
		\begin{longtable}{@{}L{0.18\textwidth}L{0.18\textwidth}L{0.18\textwidth}L{0.18\textwidth}L{0.18\textwidth}@{}}
			\toprule
			\textbf{Número de áreas}
			& \textbf{Formación con viáticos}
			& \textbf{Ensamble con viáticos}
			& \textbf{Función}
			& \textbf{Total consolidado}
			\\
			\midrule
			\endhead
			\bottomrule
			\endlastfoot
			4 áreas & \$5.440.000 COP & \$5.080.000 COP & \$8.400.000 COP & \$18.920.000 COP \\
			5 áreas & \$6.800.000 COP & \$5.080.000 COP & \$8.400.000 COP & \$20.280.000 COP \\
			6 áreas & \$8.160.000 COP & \$5.080.000 COP & \$8.400.000 COP & \$21.640.000 COP \\
			7 áreas & \$9.520.000 COP & \$5.080.000 COP & \$8.400.000 COP & \$23.000.000 COP \\
		\end{longtable}
	}
	
	\subsection{5.3 Resumen general por rangos}
	
	{\footnotesize
		\begin{longtable}{@{}L{0.18\textwidth}L{0.18\textwidth}L{0.18\textwidth}L{0.18\textwidth}L{0.18\textwidth}@{}}
			\toprule
			\textbf{Modalidad}
			& \textbf{Escala}
			& \textbf{Rango de áreas}
			& \textbf{Total mínimo estimado}
			& \textbf{Total máximo estimado}
			\\
			\midrule
			\endhead
			\bottomrule
			\endlastfoot
			Descubriendo el lenguaje & Pequeña & 4 a 7 áreas & \$12.300.000 COP & \$13.320.000 COP \\
			Explorando nuevas posibilidades & Mediana & 4 a 7 áreas & \$13.660.000 COP & \$15.700.000 COP \\
			Creación y composición intuitiva & Grande & 4 a 7 áreas & \$18.920.000 COP & \$23.000.000 COP \\
		\end{longtable}
	}
	
	\subsection{5.4 Opción adicional: cobertura audiovisual}
	
	La propuesta puede incluir una cobertura audiovisual del proceso formativo, del ensamble y de la presentación final, según los intereses de registro, memoria, comunicación o difusión de la Casa de la Cultura y de la organización de Tunjo Fest.
	
	Esta cobertura contempla el registro de cuatro sesiones de formación entre semana, realizadas en la Casa de la Cultura de Facatativá, con una duración aproximada de dos horas cada una, y una jornada final de fin de semana de hasta seis horas de disponibilidad para registrar una sesión de ensamble de aproximadamente dos horas y una presentación final de aproximadamente una hora y media. Esta disponibilidad mayor se estima porque probablemente una actividad se realice en la mañana y otra en la tarde, aunque el tiempo efectivo de grabación sea menor.
	
	La cobertura audiovisual puede contratarse en dos modalidades:
	
	{\small
		\begin{longtable}{@{}L{0.27\textwidth}L{0.30\textwidth}L{0.37\textwidth}@{}}
			\toprule
			\textbf{Modalidad de cobertura audiovisual}
			& \textbf{Alcance}
			& \textbf{Valor estimado}
			\\
			\midrule
			\endhead
			\bottomrule
			\endlastfoot
			Registro audiovisual con entrega de material en bruto & Registro de 4 sesiones entre semana de 2 horas cada una y una jornada final de hasta 6 horas durante el fin de semana. Entrega de material en bruto. & \$1.100.000 COP \\
			Registro audiovisual con entrega de material editado & Registro de 4 sesiones entre semana de 2 horas cada una y una jornada final de hasta 6 horas durante el fin de semana. Incluye edición de material audiovisual. & \$1.600.000 COP \\
		\end{longtable}
	}
	
	Estos valores pueden subir o bajar según los requerimientos de equipos, las condiciones técnicas de los espacios, el número de cámaras, las necesidades de sonido, iluminación, edición, duración final de las piezas, formatos de entrega o exigencias específicas de los eventos a registrar.
	
	La cobertura audiovisual no está incluida en el presupuesto base de las modalidades. Puede sumarse como rubro adicional, según disponibilidad presupuestal y acuerdo con la Casa de la Cultura o la organización de Tunjo Fest.
	
	\textbf{Resumen general con opciones audiovisuales}
	
	{\footnotesize
		\begin{longtable}{@{}L{0.18\textwidth}L{0.18\textwidth}L{0.18\textwidth}L{0.18\textwidth}L{0.18\textwidth}@{}}
			\toprule
			\textbf{Modalidad}
			& \textbf{Escala}
			& \textbf{Total base estimado}
			& \textbf{Con registro en bruto}
			& \textbf{Con material editado}
			\\
			\midrule
			\endhead
			\bottomrule
			\endlastfoot
			Descubriendo el lenguaje & Pequeña & \$12.300.000 COP a \$13.320.000 COP & \$13.400.000 COP a \$14.420.000 COP & \$13.900.000 COP a \$14.920.000 COP \\
			Explorando nuevas posibilidades & Mediana & \$13.660.000 COP a \$15.700.000 COP & \$14.760.000 COP a \$16.800.000 COP & \$15.260.000 COP a \$17.300.000 COP \\
			Creación y composición intuitiva & Grande & \$18.920.000 COP a \$23.000.000 COP & \$20.020.000 COP a \$24.100.000 COP & \$20.520.000 COP a \$24.600.000 COP \\
		\end{longtable}
	}
	
	\subsection{5.5 Opción adicional: artistas invitados especializados}
	
	En caso de que la Casa de la Cultura o la organización del evento decida fortalecer la visibilidad artística de la función, se puede contemplar la participación de hasta 3 artistas invitados especializados.
	
	Estos artistas tendrían honorarios propios y pueden sumarse a cualquiera de las modalidades, según disponibilidad presupuestal y pertinencia artística.
	
	{\small
		\begin{longtable}{@{}L{0.32\textwidth}L{0.62\textwidth}@{}}
			\toprule
			\textbf{Concepto}
			& \textbf{Valor}
			\\
			\midrule
			\endhead
			\bottomrule
			\endlastfoot
			1 artista invitado especial & \$600.000 COP \\
			2 artistas invitados especiales & \$1.200.000 COP \\
			3 artistas invitados especiales & \$1.800.000 COP \\
		\end{longtable}
	}
	
	\textbf{Artista invitado especial 1}
	
	Nombre: \_\_\_\_\_\_\_\_\_\_\_\_\_\_\_\_\_\_\_\_\_\_\_\_\_\_\_\_\_\_\_\_
	
	Disciplina: \_\_\_\_\_\_\_\_\_\_\_\_\_\_\_\_\_\_\_\_\_\_\_\_\_\_\_\_\_\_
	
	Breve descripción: \_\_\_\_\_\_\_\_\_\_\_\_\_\_\_\_\_\_\_\_\_\_\_\_\_\_\_\_\_\_\_\_
	
	\textbf{Artista invitado especial 2}
	
	Nombre: \_\_\_\_\_\_\_\_\_\_\_\_\_\_\_\_\_\_\_\_\_\_\_\_\_\_\_\_\_\_\_\_
	
	Disciplina: \_\_\_\_\_\_\_\_\_\_\_\_\_\_\_\_\_\_\_\_\_\_\_\_\_\_\_\_\_\_
	
	Breve descripción: \_\_\_\_\_\_\_\_\_\_\_\_\_\_\_\_\_\_\_\_\_\_\_\_\_\_\_\_\_\_\_\_
	
	\textbf{Artista invitado especial 3}
	
	Nombre: \_\_\_\_\_\_\_\_\_\_\_\_\_\_\_\_\_\_\_\_\_\_\_\_\_\_\_\_\_\_\_\_
	
	Disciplina: \_\_\_\_\_\_\_\_\_\_\_\_\_\_\_\_\_\_\_\_\_\_\_\_\_\_\_\_\_\_
	
	Breve descripción: \_\_\_\_\_\_\_\_\_\_\_\_\_\_\_\_\_\_\_\_\_\_\_\_\_\_\_\_\_\_\_\_
	
	\subsection{5.6 Cierre institucional}
	
	Tunjo Fest Soundpainting propone una experiencia artística, pedagógica y comunitaria para la Casa de la Cultura de Facatativá, mediante una práctica contemporánea, participativa e interdisciplinar.
	
	La propuesta busca dejar un resultado escénico visible, fortalecer el intercambio entre artistas invitados y artistas locales, y aportar a la programación cultural municipal.
	
	\section{Anexo presupuestal}
	
	\subsection{A.1 Costos unitarios}
	
	\emph{Todas las cifras monetarias del anexo presupuestal están expresadas en pesos colombianos COP.}
	
	{\small
		\begin{longtable}{@{}L{0.32\textwidth}L{0.62\textwidth}@{}}
			\toprule
			\textbf{Concepto}
			& \textbf{Valor unitario}
			\\
			\midrule
			\endhead
			\bottomrule
			\endlastfoot
			Hora de taller & \$150.000 COP \\
			Sesión de taller de 2 horas & \$300.000 COP \\
			Viáticos por sesión de formación & \$40.000 COP \\
			Hora de ensamble & \$1.250.000 COP \\
			Viáticos por sesión de ensamble & \$40.000 COP \\
			Intérprete en función & \$600.000 COP \\
			Invitado especial, si aplica & \$600.000 COP \\
		\end{longtable}
	}
	
	\subsection{A.2 Fórmula general de cálculo}
	
	Total = formación + viáticos de formación + ensamble + viáticos de ensamble + función + requerimientos adicionales, si aplican.
	
	Donde:
	
	\begin{itemize}
		\item
		Formación = número de áreas × número de sesiones por área × 2 horas × \$150.000 COP.
		\item
		Viáticos de formación = número de áreas × número de sesiones por área × \$40.000 COP.
		\item
		Ensamble = número de horas de ensamble × \$1.250.000 COP.
		\item
		Viáticos de ensamble = número de sesiones de ensamble × \$40.000 COP.
		\item
		Función = número de intérpretes × \$600.000 COP.
		\item
		Requerimientos adicionales = sonido, iluminación, proyección, transporte, hospedaje, alimentación, seguridad, permisos u otros rubros, según acuerdo con la Casa de la Cultura.
	\end{itemize}
	
	\subsection{A.3 Desglose de la función}
	
	La función base se calcula con 14 intérpretes:
	
	\begin{itemize}
		\item
		11 artistas del Ensamble Liminal.
		\item
		3 artistas invitados de circo de Facatativá.
	\end{itemize}
	
	Costo base de función:
	
	14 intérpretes × \$600.000 COP = \$8.400.000 COP
	
	\subsection{A.4 Desglose del ensamble por modalidad}
	
	El ensamble se desglosa en honorarios de ensamble y viáticos de ensamble. Los honorarios se calculan por hora y los viáticos se calculan por sesión.
	
	{\scriptsize
		\begin{longtable}{@{}L{0.15\textwidth}L{0.15\textwidth}L{0.15\textwidth}L{0.15\textwidth}L{0.15\textwidth}L{0.15\textwidth}@{}}
			\toprule
			\textbf{Modalidad}
			& \textbf{Horas de ensamble}
			& \textbf{Sesiones de ensamble}
			& \textbf{Honorarios de ensamble}
			& \textbf{Viáticos de ensamble}
			& \textbf{Total ensamble}
			\\
			\midrule
			\endhead
			\bottomrule
			\endlastfoot
			Descubriendo el lenguaje & 2 horas & 1 sesión & \$2.500.000 COP & \$40.000 COP & \$2.540.000 COP \\
			Explorando nuevas posibilidades & 2 horas & 1 sesión & \$2.500.000 COP & \$40.000 COP & \$2.540.000 COP \\
			Creación y composición intuitiva & 4 horas & 2 sesiones & \$5.000.000 COP & \$80.000 COP & \$5.080.000 COP \\
		\end{longtable}
	}
	
	\subsection{A.5 Desglose completo por modalidad}
	
	\textbf{Descubriendo el lenguaje}
	
	Formación: 1 sesión de 2 horas por área.
	
	Costo por área: \$300.000 COP de taller + \$40.000 COP de viáticos = \$340.000 COP.
	
	{\scriptsize
		\begin{longtable}{@{}L{0.13\textwidth}L{0.13\textwidth}L{0.13\textwidth}L{0.13\textwidth}L{0.13\textwidth}L{0.13\textwidth}L{0.13\textwidth}@{}}
			\toprule
			\textbf{Número de áreas}
			& \textbf{Formación}
			& \textbf{Viáticos de formación}
			& \textbf{Honorarios de ensamble}
			& \textbf{Viáticos de ensamble}
			& \textbf{Función}
			& \textbf{Total}
			\\
			\midrule
			\endhead
			\bottomrule
			\endlastfoot
			4 áreas & \$1.200.000 COP & \$160.000 COP & \$2.500.000 COP & \$40.000 COP & \$8.400.000 COP & \$12.300.000 COP \\
			5 áreas & \$1.500.000 COP & \$200.000 COP & \$2.500.000 COP & \$40.000 COP & \$8.400.000 COP & \$12.640.000 COP \\
			6 áreas & \$1.800.000 COP & \$240.000 COP & \$2.500.000 COP & \$40.000 COP & \$8.400.000 COP & \$12.980.000 COP \\
			7 áreas & \$2.100.000 COP & \$280.000 COP & \$2.500.000 COP & \$40.000 COP & \$8.400.000 COP & \$13.320.000 COP \\
		\end{longtable}
	}
	
	\textbf{Explorando nuevas posibilidades}
	
	Formación: 2 sesiones de 2 horas por área.
	
	Costo por área: \$600.000 COP de taller + \$80.000 COP de viáticos = \$680.000 COP.
	
	{\scriptsize
		\begin{longtable}{@{}L{0.13\textwidth}L{0.13\textwidth}L{0.13\textwidth}L{0.13\textwidth}L{0.13\textwidth}L{0.13\textwidth}L{0.13\textwidth}@{}}
			\toprule
			\textbf{Número de áreas}
			& \textbf{Formación}
			& \textbf{Viáticos de formación}
			& \textbf{Honorarios de ensamble}
			& \textbf{Viáticos de ensamble}
			& \textbf{Función}
			& \textbf{Total}
			\\
			\midrule
			\endhead
			\bottomrule
			\endlastfoot
			4 áreas & \$2.400.000 COP & \$320.000 COP & \$2.500.000 COP & \$40.000 COP & \$8.400.000 COP & \$13.660.000 COP \\
			5 áreas & \$3.000.000 COP & \$400.000 COP & \$2.500.000 COP & \$40.000 COP & \$8.400.000 COP & \$14.340.000 COP \\
			6 áreas & \$3.600.000 COP & \$480.000 COP & \$2.500.000 COP & \$40.000 COP & \$8.400.000 COP & \$15.020.000 COP \\
			7 áreas & \$4.200.000 COP & \$560.000 COP & \$2.500.000 COP & \$40.000 COP & \$8.400.000 COP & \$15.700.000 COP \\
		\end{longtable}
	}
	
	\textbf{Creación y composición intuitiva}
	
	Formación: 4 sesiones de 2 horas por área.
	
	Costo por área: \$1.200.000 COP de taller + \$160.000 COP de viáticos = \$1.360.000 COP.
	
	{\scriptsize
		\begin{longtable}{@{}L{0.13\textwidth}L{0.13\textwidth}L{0.13\textwidth}L{0.13\textwidth}L{0.13\textwidth}L{0.13\textwidth}L{0.13\textwidth}@{}}
			\toprule
			\textbf{Número de áreas}
			& \textbf{Formación}
			& \textbf{Viáticos de formación}
			& \textbf{Honorarios de ensamble}
			& \textbf{Viáticos de ensamble}
			& \textbf{Función}
			& \textbf{Total}
			\\
			\midrule
			\endhead
			\bottomrule
			\endlastfoot
			4 áreas & \$4.800.000 COP & \$640.000 COP & \$5.000.000 COP & \$80.000 COP & \$8.400.000 COP & \$18.920.000 COP \\
			5 áreas & \$6.000.000 COP & \$800.000 COP & \$5.000.000 COP & \$80.000 COP & \$8.400.000 COP & \$20.280.000 COP \\
			6 áreas & \$7.200.000 COP & \$960.000 COP & \$5.000.000 COP & \$80.000 COP & \$8.400.000 COP & \$21.640.000 COP \\
			7 áreas & \$8.400.000 COP & \$1.120.000 COP & \$5.000.000 COP & \$80.000 COP & \$8.400.000 COP & \$23.000.000 COP \\
		\end{longtable}
	}
	
	\subsection{A.6 Desglose de cobertura audiovisual}
	
	{\small
		\begin{longtable}{@{}L{0.27\textwidth}L{0.30\textwidth}L{0.37\textwidth}@{}}
			\toprule
			\textbf{Actividad a registrar}
			& \textbf{Duración estimada}
			& \textbf{Observaciones}
			\\
			\midrule
			\endhead
			\bottomrule
			\endlastfoot
			Taller entre semana 1 & 2 horas & Registro de una sesión de formación en la Casa de la Cultura de Facatativá. \\
			Taller entre semana 2 & 2 horas & Registro de una sesión de formación en la Casa de la Cultura de Facatativá. \\
			Taller entre semana 3 & 2 horas & Registro de una sesión de formación en la Casa de la Cultura de Facatativá. \\
			Taller entre semana 4 & 2 horas & Registro de una sesión de formación en la Casa de la Cultura de Facatativá. \\
			Jornada final de fin de semana & Hasta 6 horas de disponibilidad & Registro de una sesión de ensamble de aproximadamente 2 horas y de una presentación final de aproximadamente 1 hora y media. Se estima una disponibilidad mayor porque las actividades podrían realizarse en horarios separados, una en la mañana y otra en la tarde. \\
		\end{longtable}
	}
	
	\textbf{Modalidades de entrega:}
	
	{\small
		\begin{longtable}{@{}L{0.27\textwidth}L{0.30\textwidth}L{0.37\textwidth}@{}}
			\toprule
			\textbf{Modalidad}
			& \textbf{Valor estimado}
			& \textbf{Entrega}
			\\
			\midrule
			\endhead
			\bottomrule
			\endlastfoot
			Registro con material en bruto & \$1.100.000 COP & Entrega del material grabado en bruto. \\
			Registro con material editado & \$1.600.000 COP & Entrega de material audiovisual editado. \\
		\end{longtable}
	}
	
	\subsection{A.7 Anexos recomendados}
	
	\begin{itemize}
		\item
		Portafolio de Soundpainting Colombia.
		\item
		Portafolio del Ensamble Liminal.
		\item
		Hojas de vida breves del equipo artístico.
		\item
		Hojas de vida breves de los 3 artistas invitados de circo.
		\item
		Hojas de vida breves de los invitados especiales, si aplica.
		\item
		Rider técnico.
		\item
		Presupuesto por modalidad.
		\item
		Cronograma detallado.
		\item
		Video de referencia.
		\item
		Enlace de presentación en Canva.
		\item
		Registro de experiencias previas.
		\item
		Plan básico de riesgos.
	\end{itemize}
	
\end{document}
